
\section{Conclusion}
We proposed ProtoDiff, a novel approach to prototype-based meta-learning that utilizes a task-guided diffusion model during the meta-training phase to generate efficient class representations. We addressed the limitations of estimating a deterministic prototype from a limited number of examples by optimizing a set of prototypes to accurately represent individual tasks and training a task-guided diffusion process to model each task's underlying distribution. Our approach considers both probabilistic and task-guided prototypes, enabling efficient adaptation to new tasks while maintaining the informativeness and scalability of prototypical networks. Extensive experiments on three distinct few-shot learning scenarios: within-domain, cross-domain, and few-task few-shot learning validated the effectiveness of ProtoDiff. Our results demonstrated significant improvements in classification accuracy compared to state-of-the-art meta-learning techniques, highlighting the potential of task-guided diffusion in augmenting few-shot learning and advancing meta-learning.


\section*{Acknowledgment}
This work is financially supported by the Inception Institute of Artificial Intelligence, the University of Amsterdam and the allowance 
Top consortia for Knowledge and Innovation (TKIs) from the Netherlands Ministry of Economic Affairs and Climate Policy.